\corrigesChapitre{6}{Reconnaître une transformation chimique}
\par\noindent\begin{minipage}{\linewidth}
\begin{corrige}{1}{Un changement physique}
\begin{questions}
\item Non : les molécules restent celles de l’eau.
\item C’est une transformation physique, un changement d’état.
\end{questions}
\end{corrige}
\end{minipage}
\par\noindent\begin{minipage}{\linewidth}
\begin{corrige}{2}{Réactif ou produit}
\begin{questions}
\item Un réactif.
\item Un produit.
\end{questions}
\end{corrige}
\end{minipage}
\par\noindent\begin{minipage}{\linewidth}
\begin{corrige}{3}{Le trouble observé}
\begin{questions}
\item Du dioxyde de carbone.
\item Non : il s’agit ici d’un test d’identification qualitatif.
\end{questions}
\end{corrige}
\end{minipage}
\par\noindent\begin{minipage}{\linewidth}
\begin{corrige}{4}{Le fer rouille}
\begin{questions}
\item Non : c’est une transformation chimique.
\item Le fer et d’autres réactifs sont consommés pour former notamment des espèces constituant la rouille.
\end{questions}
\end{corrige}
\end{minipage}
\par\noindent\begin{minipage}{\linewidth}
\begin{corrige}{5}{Les bulles ne suffisent pas}
\begin{questions}
\item Non.
\item La vaporisation crée des bulles de vapeur d’eau ; il n’est pas nécessaire de supposer une nouvelle espèce chimique.
\end{questions}
\end{corrige}
\end{minipage}
\par\noindent\begin{minipage}{\linewidth}
\begin{corrige}{6}{Lire un bilan}
\begin{questions}
\item Réactifs : carbone et dioxygène ; produit : dioxyde de carbone.
\item Deux atomes O de chaque côté ; un C de chaque côté.
\end{questions}
\end{corrige}
\end{minipage}
\par\noindent\begin{minipage}{\linewidth}
\begin{corrige}{7}{Un coefficient utile}
\begin{questions}
\item Avant : quatre H et deux O ; après : quatre H et deux O.
\item Deux molécules d’eau, dans ce bilan de comptage.
\end{questions}
\pcFigureEx{figures/cor-reaction-eau.pdf}{Les coefficients multiplient des molécules entières ; les indices des espèces restent inchangés.}
\end{corrige}
\end{minipage}
\par\noindent\begin{minipage}{\linewidth}
\begin{corrige}{8}{Le témoin}
\begin{questions}
\item Il sert de témoin pour vérifier l’état initial et l’effet de l’ajout.
\item Le protocole n’est pas interprétable tel quel : réactif ou matériel peuvent déjà être contaminés.
\end{questions}
\pcFigureEx{figures/cor-temoin.pdf}{Le témoin recherche une cause de trouble indépendante de l’ajout de gaz.}
\end{corrige}
\end{minipage}
\par\noindent\begin{minipage}{\linewidth}
\begin{corrige}{9}{Une identification limitée}
\begin{questions}
\item Non : le test peut révéler une espèce au sein d’un mélange.
\item Le prélèvement contient du dioxyde de carbone détectable dans les conditions du test.
\end{questions}
\end{corrige}
\end{minipage}
\par\noindent\begin{minipage}{\linewidth}
\begin{corrige}{10}{Un test négatif}
\begin{questions}
\item Non : une quantité peut être trop faible ou le test mal réalisé.
\item Vérifier le réactif avec un témoin positif et vérifier les conditions de prélèvement et de contact.
\end{questions}
\end{corrige}
\end{minipage}
\par\noindent\begin{minipage}{\linewidth}
\begin{corrige}{11}{Corriger un équilibrage}
\begin{questions}
\item Il change l’espèce produite au lieu de conserver sa formule.
\item $2\,\ce{H2}+\ce{O2}\rightarrow2\,\ce{H2O}$ conserve les atomes et forme bien de l’eau.
\end{questions}
\par\smallskip\noindent\textbf{Pour comprendre.}\enspace Changer H2O en H2O2 change le produit : on obtient une autre espèce. Il faut modifier les coefficients : deux H2 et un O2 donnent deux H2O, soit quatre H et deux O de chaque côté.\par
\pcFigureEx{figures/cor-reaction-eau.pdf}{Les coefficients multiplient des molécules entières ; les indices des espèces restent inchangés.}
\end{corrige}
\end{minipage}
\par\noindent\begin{minipage}{\linewidth}
\begin{corrige}{12}{Séparer les faits des idées}
\begin{questions}
\item La première décrit un fait observable ; la seconde est une interprétation chimique.
\item Identifier le gaz avec un test adapté et contrôler qu’il n’était pas simplement dissous avant la transformation.
\end{questions}
\end{corrige}
\end{minipage}
\par\noindent\begin{minipage}{\linewidth}
\begin{corrige}{13}{Un métal et un acide}
\begin{questions}
\item Les réactivités dépendent du métal et de la solution ; la banque ne décrit que la situation annoncée.
\item Jamais. On utilise une méthode autorisée de mesure de pH, avec les protections prescrites.
\end{questions}
\end{corrige}
\end{minipage}
\par\noindent\begin{minipage}{\linewidth}
\begin{corrige}{14}{Un gaz déjà présent}
\begin{questions}
\item Non : il identifie le gaz, pas son origine.
\item Étudier les espèces présentes avant et après, ainsi que les conditions de l’expérience.
\end{questions}
\end{corrige}
\end{minipage}
\par\noindent\begin{minipage}{\linewidth}
\begin{corrige}{15}{Concevoir une preuve}
\begin{questions}
\item Chercher une nouvelle espèce par un test adapté et vérifier la consommation d’un réactif grâce à des informations complémentaires.
\item Une dissolution peut rendre la poudre invisible sans constituer à elle seule une réaction au sens du bilan recherché.
\end{questions}
\end{corrige}
\end{minipage}
\par\noindent\begin{minipage}{\linewidth}
\phantomsection\label{cor-defi-6}\pcSolutionDefi{6}{Deux flacons effervescents}\begin{questions}
\item Ils montrent la présence de dioxyde de carbone détectable dans les deux prélèvements.
\item Dans A, le gaz était déjà dissous ; dans B, la banque et les observations avant-après indiquent sa formation chimique.
\item Le témoin vérifie que l’eau de chaux ne trouble pas spontanément. Même gaz identifié ne signifie pas même transformation.
\end{questions}
\end{minipage}
\par\noindent\begin{minipage}{\linewidth}
\phantomsection\label{cor-activite-6}\pcSousTitre{Pistes pour l’investigation}
Les deux prélèvements peuvent contenir du dioxyde de carbone et troubler l’eau de chaux. Il faut ensuite examiner la situation initiale : le dégazage sépare un gaz dissous, tandis que la réaction carbonate-acide forme notamment du gaz à partir des réactifs. Le témoin contrôle le comportement de l’eau de chaux sans ajout. Les seules bulles ne distinguent pas ces mécanismes.\par\end{minipage}
